% countdown-srs.tex
% Software Requirements Specification for the countdown program.
% Structure per IEEE Std 830-1998 / ISO/IEC/IEEE 29148:2018.
% Compiled with: pdflatex countdown-srs.tex

\documentclass[12pt]{article}

\usepackage{geometry}
\geometry{letterpaper, margin=1in}
\usepackage{hyperref}
\usepackage{booktabs}
\usepackage{array}
\usepackage{longtable}
\usepackage{enumitem}
\usepackage{titlesec}
\usepackage{fancyhdr}
\usepackage{xcolor}

\definecolor{cdteal}{RGB}{0,90,120}

\pagestyle{fancy}
\fancyhf{}
\renewcommand{\headrulewidth}{0.4pt}
\fancyhead[L]{\small Countdown Project}
\fancyhead[R]{\small SRS: countdown}
\fancyfoot[C]{\thepage}

\titleformat{\section}{\large\bfseries\color{cdteal}}{\thesection}{0.5em}{}[\titlerule]
\titleformat{\subsection}{\normalsize\bfseries}{\thesubsection}{0.5em}{}

\newcommand{\code}[1]{\texttt{#1}}
\newcommand{\req}[1]{\texttt{#1}}

% A requirement table row: id | statement
\newcommand{\rrow}[2]{\req{#1} & #2 \\}

\begin{document}

\begin{titlepage}
  \centering
  \vspace*{2cm}
  {\Huge\bfseries\color{cdteal} Software Requirements\\ Specification\\[0.5em]}
  {\Large countdown: An X~Window Countdown Timer}\\[2em]
  {\large Countdown Software Project}\\[0.5em]
  {\normalsize Prepared by: J.\,R.\,Evans}\\[2em]
  \begin{tabular}{ll}
    Document identifier: & CD-SRS-001                                     \\
    Revision:            & 1.1                                            \\
    Date:                & June 15, 2026                                  \\
    Status:              & Released                                       \\
    Governing standard:  & IEEE Std 830-1998 (SRS)                        \\
  \end{tabular}
  \vfill
  {\small\textit{This document is prepared in accordance with the
  principles of rational documentation articulated in
  D.\,L.\,Parnas and P.\,C.\,Clements, ``A Rational Design Process:
  How and Why to Fake It,'' \textit{IEEE Transactions on Software
  Engineering}, vol.\,SE-12, no.\,2, pp.\,251--257, February 1986.}}
\end{titlepage}

\tableofcontents
\clearpage

% =================================================================
\section{Introduction}
% =================================================================

\subsection{Purpose}

This Software Requirements Specification (SRS) states the externally
observable requirements on the \code{countdown} program.  It is the
contract against which the Software Design Document (CD-SDD-001) is
written and against which the Software Test Plan (CD-STP-001) and its
automated harness verify the delivered program.

\subsection{Scope}

\code{countdown} is an X~Window System application that displays the time
remaining until a fixed future instant and updates that display live.  It
runs in two modes: by default it counts down in the controlling
\emph{terminal}; given the \code{-gui} option it instead opens a single
X~window.  The SRS covers the program tangled from \code{countdown.w}.  It
does not cover the X~server, the operating system's time-zone database, or
the CWEB tools themselves, except to state the program's reliance on them.

\subsection{Definitions and acronyms}

\begin{longtable}{p{0.22\textwidth}p{0.72\textwidth}}
\toprule
\textbf{Term} & \textbf{Meaning} \\
\midrule
Target instant & The fixed future date and time to which the program counts down. \\
Armed & Started but not yet running; shows the full remaining duration. \\
Running & Counting down live. \\
Paused & Display frozen at the value captured when the user paused. \\
Arrived & The remaining time has reached zero. \\
Terminal mode & The default mode: the countdown runs in the controlling terminal, no window. \\
GUI mode & Selected by \code{-gui}: the countdown runs in an X~window. \\
CWEB & Knuth's literate-programming system; \code{ctangle} extracts C, \code{cweave} extracts \TeX. \\
Xlib & The C-language client interface to the X~Window System. \\
POSIX & IEEE Std 1003.1 operating-system interface. \\
\bottomrule
\end{longtable}

\subsection{Rational-process note}

In the spirit of Parnas and Clements, the requirements below are presented
as a clean, numbered hierarchy as though derived top-down.  They were in
fact refined alongside the design; the idealised presentation is retained
because it is the form most useful to a reader.

\subsection{Requirement identifier scheme}

Each requirement carries an identifier \req{REQ-\textit{C}-\textit{nn}}
where \textit{C} is a category letter---\textbf{F}unctional,
\textbf{G}UI, \textbf{T}iming, \textbf{D}esign-constraint, or
\textbf{B}uild---and \textit{nn} is a serial number.  Identifiers are
stable and are referenced verbatim by CD-STP-001 and by the automated test
harness.

% =================================================================
\section{Overall description}
% =================================================================

\subsection{Product perspective}

\code{countdown} is a self-contained, standalone program.  It depends on
an X~server for display and input, and on the POSIX time facilities of the
host for the current time and time-zone conversion.  It has no persistent
storage, no configuration file, and no network interface.

\subsection{Product functions (summary)}

The program converts a target calendar string to an absolute instant,
repeatedly computes the seconds remaining until that instant, formats them
as days and a clock, and paints them in an X~window, reacting to keyboard
and mouse input to start, pause, resume, and quit.

\subsection{User characteristics}

The user is anyone able to launch a program in an X~session.  No
specialised knowledge is assumed beyond reading a window and pressing a
key or mouse button.

\subsection{Constraints}

The program shall be written as a literate program in CWEB and built with
the project's \code{cweb.sh} helper.  It shall link only against core
Xlib.  It shall run on Linux/X11 and macOS/XQuartz hosts.

\subsection{Assumptions and dependencies}

It is assumed that an X~server is reachable via the \code{DISPLAY}
environment variable, that the host clock is set correctly, and that the
time-zone database reflects the rules in force for the target date.

% =================================================================
\section{Specific requirements}
% =================================================================

\subsection{Functional requirements}

\begin{longtable}{p{0.16\textwidth}p{0.78\textwidth}}
\toprule
\textbf{ID} & \textbf{Requirement} \\
\midrule
\endhead
\rrow{REQ-F-01}{The program shall count down to a \emph{target instant}
  specified as local calendar time.}
\rrow{REQ-F-02}{In the absence of a command-line argument, the target
  instant shall default to \code{2026-08-19 08:00:00} local time.}
\rrow{REQ-F-03}{The program shall accept an alternate target instant as
  its first command-line argument, in the form
  \code{"YYYY-MM-DD HH:MM:SS"}.}
\rrow{REQ-F-04}{Given a target argument that is malformed or denotes an
  invalid date, the program shall print a diagnostic to standard error
  and terminate with a non-zero exit status, without opening a window.}
\rrow{REQ-F-05}{The program shall display the remaining time as a whole
  number of days followed by a 24-hour \code{HH:MM:SS} clock, with minutes
  and seconds zero-padded.}
\rrow{REQ-F-06}{When the target instant has been reached, the remaining
  time shall read \code{0 d 00:00:00}; the program shall never display a
  negative remaining time.}
\rrow{REQ-F-07}{Upon arrival the program shall display a status message
  announcing that the target instant has been reached.}
\rrow{REQ-F-08}{The program shall run in the controlling terminal by
  default, and shall open an X~window instead when given the \code{-gui}
  option as a command-line argument.  The \code{-gui} option and the
  optional target argument may appear in any order.}
\rrow{REQ-F-09}{In terminal mode the program shall update the remaining
  time on a single line, rewriting it in place once per second (without
  scrolling), and on arrival append the arrival message; it shall require
  no X~server.}
\bottomrule
\end{longtable}

\subsection{Graphical user interface requirements}

\begin{longtable}{p{0.16\textwidth}p{0.78\textwidth}}
\toprule
\textbf{ID} & \textbf{Requirement} \\
\midrule
\endhead
\rrow{REQ-G-01}{In \code{-gui} mode the program shall present its output in
  a graphical window on the X~Window System using the Xlib client
  interface.}
\rrow{REQ-G-02}{The window shall display, simultaneously, the target
  instant in human-readable form, the remaining time, a status line, and a
  one-line hint naming the available controls.}
\rrow{REQ-G-03}{The user shall be able to start, pause, and resume the
  countdown by pressing the space bar or by clicking a mouse button in the
  window.}
\rrow{REQ-G-04}{The user shall be able to terminate the program by pressing
  \code{q} or the Escape key.}
\rrow{REQ-G-05}{The window shall be repainted in response to expose and
  resize (configure) events so that its contents remain correct.}
\rrow{REQ-G-06}{If the preferred font is unavailable, the program shall
  fall back to the server's default font rather than fail.}
\bottomrule
\end{longtable}

\subsection{Timing requirements}

\begin{longtable}{p{0.16\textwidth}p{0.78\textwidth}}
\toprule
\textbf{ID} & \textbf{Requirement} \\
\midrule
\endhead
\rrow{REQ-T-01}{While running, the displayed remaining time shall advance
  at least once per second.}
\rrow{REQ-T-02}{The remaining time shall be computed from the host's
  wall-clock time as obtained from the POSIX \code{time} service.}
\rrow{REQ-T-03}{Conversion of the target calendar string to an absolute
  instant shall honour the host's local time zone and daylight-saving
  rules, as provided by the POSIX \code{mktime} service.}
\rrow{REQ-T-04}{The program shall remain idle between updates by blocking
  on the POSIX \code{select} service with a timeout; it shall not
  busy-wait.}
\bottomrule
\end{longtable}

\subsection{Design-constraint requirements}

\begin{longtable}{p{0.16\textwidth}p{0.78\textwidth}}
\toprule
\textbf{ID} & \textbf{Requirement} \\
\midrule
\endhead
\rrow{REQ-D-01}{The program shall be written as a literate program in the
  CWEB system, with \code{countdown.w} as the single source of record.}
\rrow{REQ-D-02}{Each woven module (CWEB section) shall contain fewer than
  twenty-four lines of code.}
\rrow{REQ-D-03}{Every POSIX system call used by the program (directly or
  on the program's behalf by Xlib) shall be documented in a dedicated
  section of the literate program and entered in its index.}
\rrow{REQ-D-04}{The compilable C source shall be a build artefact
  regenerable in full from \code{countdown.w} by \code{ctangle}.}
\rrow{REQ-D-05}{The tangled C source shall compile cleanly under
  \code{-Wall} with no warnings.}
\bottomrule
\end{longtable}

\subsection{Build and portability requirements}

\begin{longtable}{p{0.16\textwidth}p{0.78\textwidth}}
\toprule
\textbf{ID} & \textbf{Requirement} \\
\midrule
\endhead
\rrow{REQ-B-01}{The program shall be tangled, woven, and compiled through
  the project \code{Makefile}, which drives the CWEB tools via the
  \code{cweb.sh} helper.}
\rrow{REQ-B-02}{The program shall link against the core Xlib library only,
  with no dependency on any widget toolkit.}
\rrow{REQ-B-03}{Before any container-driven step, the \code{Makefile} shall
  invoke \code{docker-check.sh} to verify that the Docker daemon is
  running and start it (Docker Desktop) if it is not.}
\rrow{REQ-B-04}{The \code{Makefile} shall build all documentation---the
  woven program listing and the IEEE/ISO document set---inside the same
  container used by \code{cweb.sh} (\code{make docs}), and shall provide a
  \code{clean} target that removes all intermediate build files.}
\rrow{REQ-B-05}{The \code{Makefile} shall install the final executable into
  \code{/usr/local/bin} (\code{make install}).}
\bottomrule
\end{longtable}

\subsection{External interface requirements}

\textbf{Command-line interface.}  \code{countdown [-gui] [TARGET]}, where
the optional \code{-gui} option selects the X~window mode (the default is a
terminal countdown, REQ-F-08) and the optional \code{TARGET} is a string
of the form \code{"YYYY-MM-DD HH:MM:SS"} (REQ-F-03); the two may appear in
either order.  Exit status is \code{0} on normal termination, \code{1} on
inability to open the display (\code{-gui} mode only), and \code{2} on a
malformed target (REQ-F-04).

\medskip\noindent
\textbf{User interface.}  One top-level X~window (REQ-G-01--REQ-G-05) with
keyboard and pointer input.

\medskip\noindent
\textbf{Environment interface.}  The \code{DISPLAY} variable selects the
X~server; the \code{TZ} variable (or system default) selects the time
zone used by \code{mktime}.

\subsection{Non-functional requirements}

\textbf{Performance.}  Idle CPU use shall be negligible, achieved by
\code{select}-based blocking (REQ-T-04).  \textbf{Portability.}  The
program shall build and run on Linux/X11 and macOS/XQuartz (Section~2.4).
\textbf{Readability.}  The program shall be legible as a typeset document
produced by \code{cweave} (REQ-D-01, REQ-D-02).

% =================================================================
\section{Requirements traceability}
% =================================================================

Each requirement traces upward to a ConOps capability or scenario
(CD-CONOPS-001) and downward to one or more design modules
(CD-SDD-001) and test cases (CD-STP-001).  The condensed upward trace is:

\begin{longtable}{p{0.20\textwidth}p{0.74\textwidth}}
\toprule
\textbf{Requirement(s)} & \textbf{ConOps origin} \\
\midrule
REQ-F-01, REQ-F-02 & \S6.3 capabilities 1; \S7.1 Scenario~1 \\
REQ-F-03 & \S7.2 Scenario~2 (custom target) \\
REQ-F-04 & \S7.4 Scenario~4 (mistyped target) \\
REQ-F-05--REQ-F-07 & \S6.3 capabilities 2 and 4; \S7.3 Scenario~3 \\
REQ-F-08, REQ-F-09 & \S6.4 modes of operation; \S7.5 Scenario~5 (terminal) \\
REQ-G-01--REQ-G-06 & \S6.1 overview; \S6.4 modes of operation \\
REQ-T-01--REQ-T-04 & \S6.2 operational environment; \S8.2 trade-offs \\
REQ-D-01--REQ-D-05 & \S4 justification (transparency); \S8 analysis \\
REQ-B-01--REQ-B-05 & \S6.2 operational environment; build/packaging \\
\bottomrule
\end{longtable}

\end{document}
